\clearpage
\fancyhead[L]{\small COMET.}
\SourceLeaf{425}
\begin{center}
\Large COMET.

\medskip
Göttingische gelehrte Anzeigen. Stück 83. Seite 825..829. 1819 Mai 24.
\end{center}

\par\indent Im 28. Stück dieser Anzeigen [S. 417 d. B.] haben wir eine Nachricht über die beiden
Cometen vom vorigen Jahre mitgetheilt, und dabei bereits die Wahrscheinlichkeit der
Identität des einen dieses Cometen mit dem ersten von 1805 angedeutet. Diese
Wahrscheinlichkeit ist durch die spätern Rechnungen des Hrn. Prof. \textsc{Encke} so sehr
erhöhet worden, dass kaum noch ein Zweifel an der Identität statt finden kann. Wir
tragen kein Bedenken, diese Bereicherung der Kenntniss unsers Sonnensystems zu den
merkwürdigsten Entdeckungen des neunzehnten Jahrhunderts zu zählen, und theilen hier
darüber das Hauptsächlichste, aus einigen Briefen jenes Astronomen an Hrn Hofr.
\textsc{Gauss}, mit. Da noch immer von dem Cometen von 1818 keine weitern brauchbaren
Beobachtungen, als die im angeführten Blatte
abgedruckten, bekannt geworden sind, so hielt Hr. Prof. \textsc{Encke} vorerst für das
Rathsamste, eine neue Discussion der Beobachtungen des Cometen von 1805 vorzunehmen.
Es zeigt sich hier das überraschende Resultat, dass eine Ellipse mit einer
Umlaufszeit von 3.42 Jahren die Beobachtungen am besten darstelle; für den Cometen
von 1818 hatte er nach den angeführten Elementen eine Umlaufszeit von 3.59 Jahren
gefunden: beide Angaben involvirten natürlich noch einige Unzuverlässigkeit. Da nun
zwischen den beiden beobachteten Durchgangszeiten durch die Sonnennähe 13.19 Jahre
verflossen sind, so war unter der Voraussetzung der Identität das Wahrscheinlichste,
dass der Comet in der Zwischenzeit vier Umlaufe gemacht habe (Dass der, wenn auch
nicht zu den allerschwächsten gehörende, doch nur telescopische Comet, bei seiner
dreimaligen Wiederkunft zur Sonnennähe 1809, 1812 und 1815 unbemerkt geblieben ist,
kann keinen Unterrichteten befremden). Mit einer diesem Schlusse gemäss abgeänderten
Umlaufszeit berechnete Hr. Prof. \textsc{Encke} neue elliptische Elemente nach den
Beobachtungen von 1805, wodurch diese auf das befriedigendste dargestellt wurden.
Ueberraschend angenehm war hiebei, dass diese neue Bahn in den einzelnen Elementen
der aus den Beobachtungen von 1818 berechneten ungemein viel näher gekommen war.
Hr. \textsc{Encke} entschloss sich, nach diesem aufmunternden Erfolge, die Störungen, welche
jene Elemente während der ganzen 13 Jahre durch die Einwirkung des Jupiter erleiden
mussten, nach der Methode des Hrn. Hofrath \textsc{Gauss} zu berechnen. Auch diese mühsame
Arbeit blieb nicht unbelohnt: eine noch bessere Uebereinstimmung war die Folge davon.
Wir setzen hier diese Elemente her, wie sie also, bloss auf die Beobachtungen von
1805 gegründet, und nur in Rücksicht der
grossen Axe so bestimmt, dass nach vier Umläufen die Durchgangszeit durch die
Sonnennähe mit der 1819 beobachteten zusammenfiel, nach den Einwirkungen des Jupiter
für den Anfang des Jahrs 1819 durch Hrn. \textsc{Encke}'s Rechnung gefunden worden sind:

\begin{center}
\begin{tabular}{l l}
\multicolumn{2}{l}{Durchgang durch die Sonnennähe 1819 Jan. 27.27} \\
Länge der Sonnennähe . . . . . . . . . . . . & $156^\circ$ $59'$ $30''$ \\
Länge des aufsteigenden Knoten . . . . . . . . . & 334\ \ 31\ \ \ 0 \\
\multicolumn{2}{l}{\quad beide vom mittlern Aequinoctium von 1819 gezählt} \\
Neigung der Bahn . . . . . . . . . . . . . . & 13\ \ 36\ \ 35 \\
Excentricitätswinkel . . . . . . . . . . . . . & 58\ \ \ 2\ \ 58 \\
Logarithm der halben grossen Axe . . . . . . . . & 0.3451979 \\
\end{tabular}
\end{center}

\Needspace{12\baselineskip}\par\indent Endlich machte Hr. \textsc{Encke} noch eine neue Bestimmung der Bahn blos aus den Beobachtungen
\SourceLeaf{426}von 1818..1819, indem er den Logarithm der halben grossen Axe $= 0.345$ voraussetzte, und fand, dass
damit die Beobachtungen sich noch viel besser darstellen liessen, als in der früher berechneten Bahn.
Diese neuen Elemente sind:

\begin{center}
\begin{tabular}{l l}
Zeit der Sonnennähe 1819 Januar . . . . . . . . & 27.27545 Seeberger Zeit \\
Länge der Sonnennähe . . . . . . . . . . . . & $157^\circ\,5'\,53$: \\ 
Länge des Knoten . . . . . . . . . . . . . . & $334\;\;43\;\;37$ \\
Neigung der Bahn . . . . . . . . . . . . . . & $13\;\;38\;\;42$ \\
Excentricitätswinkel . . . . . . . . . . . . . & $58\;\;6\;\;45.52$ \\
\end{tabular}
\end{center}

\par\indent Die Uebereinstimmung dieser Elemente mit den aus den Beobachtungen von 1805 abgeleiteten ist
so gross, dass die Wahrscheinlichkeit der Identität beider Cometen derjenigen, die man bei nicht rein
mathematischen Gegenständen Gewissheit nennt, gleich kommt.

\par\indent Wünschen könnte man nun noch, mit Bestimmtheit zu sehen, in wie fern diese herrliche Uebereinstimmung
von der vorausgesetzten Anzahl von \emph{vier} Umläufen abhängig ist, und es scheint uns von Wichtigkeit zu
sein, dass auch noch die Bahnen bestimmt werden, wobei nur drei Umläufe vorausgesetzt werden.
Voraussehen lässt sich, dass sowohl die Beobachtungen von 1805, als die von 1818 durch Ellipsen von
4.4 Jahren Umlaufszeit auch noch recht gut sich werden darstellen lassen; allein nach einem vorläufigen
Ueberschlage des Hrn. Hofr. \textsc{Gauss} wird die Uebereinstimmung beider Bahnen bedeutend geringer sein,
als die der oben angeführten. Wir hoffen, dass Hr. \textsc{Encke}, der sich durch die bisherigen Rechnungen
schon ein so grosses Verdienst erworben hat, sich auch dieser Untersuchung unterziehen werde, und wir
sind geneigt zu glauben. 
dass das Resultat so ausfallen werde, dass man auch ohne die Störungen
vollständig zu berechnen über die Unstatthaftigkeit der Voraussetzung einer andern Anzahl von Umläufen
als vier, werde urtheilen können. Unnöthig würde übrigens diese Discussion werden, wenn man in diesem
Jahre den Cometen noch einmal beobachten könnte. In unsern nördlichen Gegenden ist dazu freilich wenig
oder gar keine Hoffnung; allein in Italien, bei Anwendung sehr lichtstarker Instrumente, scheint es doch
nicht ganz unmöglich zu sein. Wir theilen daher zur Erleichterung solcher Versuche hier noch eine von
Hrn. \textsc{Encke} berechnete Ephemeride für die nächsten Monate mit. Die Zeiten sind $4^{\mathrm{u}}\,26^{\mathrm{m}}$ Minuten
Seeberger Zeit.
\nopagebreak[4]
\begin{center}
\begin{tabular}{c|c|c|c}
1819 & Gerade & Südliche & L. der Entf. \\
 & Aufsteigung & Abweichung & von d. Erde \\
\hline
Mai\quad 31 & $336^\circ\,28'$ & $19^\circ\,43'$ & 0.2277 \\
Juni\quad 10 & $335\;\;27$ & $20\;\;22$ & 0.2168 \\
\phantom{Juni\quad} 20 & $333\;\;30$ & $21\;\;19$ & 0.2070 \\
\phantom{Juni\quad} 30 & $330\;\;36$ & $22\;\;28$ & 0.2004 \\
Juli\quad 10 & $326\;\;49$ & $23\;\;43$ & 0.1991 \\
\phantom{Juli\quad} 20 & $322\;\;26$ & $24\;\;54$ & 0.2049 \\
\phantom{Juli\quad} 30 & $317\;\;47$ & $25\;\;51$ & 0.2189 \\
Aug.\quad 9 & $313\;\;17$ & $26\;\;30$ & 0.2408 \\
\phantom{Aug.\quad} 19 & $309\;\;20$ & $26\;\;47$ & 0.2691 \\
\phantom{Aug.\quad} 29 & $306\;\;11$ & $26\;\;48$ & 0.3020 \\
Sept.\quad 8 & $303\;\;54$ & $26\;\;35$ & 0.3374 \\
\phantom{Sept.\quad} 18 & $302\;\;29$ & $26\;\;13$ & 0.3737 \\
\end{tabular}
\end{center}

\Needspace{7\baselineskip}\par\indent Dabei ist noch zu bemerken, dass diese Ephemeride nach den frühern Elementen berechnet ist.
Eine Vorstellung, wie viel die neuern davon abweichen, geben folgende zwei nach diesen berechnete Oerter:

\begin{center}
\begin{tabular}{c|c|c|c}
Juli\quad 30 & $320^\circ\,32$ & $24^\circ\,36'$ & 0.2097 \\
Aug.\quad\;\;9 & $315\;\;59$ & $25\;\;21$ & 0.2294 \\
\end{tabular}
\end{center}

\SourceLeaf{427}\par\indent Wir schliessen diese Anzeige mit dem Wunsche, dass die höchst merkwürdige Entdeckung die
Freunde der Astronomie zu einem verdoppelten Eifer in Aufsuchung von Cometen anfeuern möge: denn
alles lässt vermuthen, dass jene nur der Anfang zu einer unermesslichen und nach und nach reifenden
Erndte sein wird.
